\documentclass[runningheads]{llncs}

\usepackage[T5]{fontenc}
\usepackage[utf8]{inputenc}
\usepackage{booktabs}
\usepackage{amsmath}
\usepackage{hyperref}
\usepackage{xcolor}

\newcommand{\todoitem}[1]{\par\noindent\textcolor{red}{[TODO: #1]}\par}

\title{Early Course-Quality Prediction under a Small Training Budget}
\titlerunning{Early Course-Quality Prediction}
\author{Anonymous Authors}
\authorrunning{Anonymous Authors}
\institute{Anonymous Institution}

\begin{document}
\maketitle

\begin{abstract}
We study training-set compression for early course-quality prediction from four
cumulative behavioural snapshots.  The development protocol evaluates both
enrolment-level utility (UL) and course-level estimation error (EC), while
measuring selector time and memory under a fixed training budget.  This M0
manuscript records the data and evaluation protocol only.  Candidate selection,
confirmatory test results, and performance claims remain intentionally sealed
until the M1 freeze.
\keywords{Course quality prediction \and coreset selection \and early prediction \and imbalanced learning}
\end{abstract}

\section{Introduction}
Early warning systems are most useful when they can be updated repeatedly as a
course unfolds.  In our setting, each learner--course pair is represented at
25\%, 50\%, 75\%, and 90\% of its observation window.  Training on all four
snapshots is expensive and highly redundant, which motivates selecting a small,
weighted subset of the development training pool.

Selection quality cannot be judged by enrolment classification alone.  The
same predictions are aggregated to estimate course-level quality, and a method
may improve macro-F1 while distorting class mass or the resulting course
estimate.  We therefore register two endpoints: UL at the enrolment level and
EC at the course level.  Rare-class metrics, calibration, exact or realised
budget, wall time, and memory are reported rather than hidden behind one score.

At M0, our contribution is a leakage-controlled experimental protocol and an
auditable implementation.  It provides course-disjoint development data,
train-only preprocessing, four-phase reuse of one fitted model, baseline
fidelity checks, and measured resource gates.  The final methodological
contribution is deliberately not stated before candidate selection and the
novelty audit are complete.

\section{Related Work}
\subsection{Early Course-Quality Prediction}
Prior early-prediction systems motivate phase-specific behavioural snapshots
and the separation of learner-level prediction from course-level monitoring.
Our operational target is a formula-derived quality index rather than a
human-validated judgement, so all claims are restricted to this definition and
dataset.  A direct feature-by-feature comparison with TEMPO-DQ and LoDiBi is
retained as an M0 approval item rather than asserted from names alone.
\subsection{Coresets and Data Selection for Tabular Classification}
Dataset condensation can construct synthetic training examples by matching
training gradients \cite{zhao2021datasetcondensation}; this differs from our
subset-selection contract, which must return traceable source rows.  Kernel
thinning provides distribution-compression guarantees in an RKHS but has a
quadratic-time form in its general construction \cite{dwivedi2024kernel}; our
classwise implementation is therefore resource-gated and labelled as an
adaptation.  CoreTab is included from a commit-pinned software source
\cite{hadarcoretabsoftware}; its native interface is binary, and the multiclass
one-vs-all baseline is explicitly labelled adapted rather than native.
\subsection{Multiclass and Imbalanced Selection}
The three classes are strongly imbalanced, making overall accuracy inadequate.
Our controls include stratified random selection and phase--class stratified
selection.  Candidate methods must preserve feasible class coverage and report
realised mass.  The one-vs-all CoreTab adapter allocates class quotas by largest
remainder and never pads missing rows.  These engineering choices are controls,
not novelty claims.

\section{Problem Formulation}
Let $k$ denote a learner, $c$ a course, and
$p\in\{1,2,3,4\}$ an observation phase.  The vector $x_{kcp}$ contains only
information observed no later than the phase cutoff.  The final quality index
$q_{kc}\in[0,1]$ is fixed across phases and lower values are better.  Its class
label is
\[
y_{kc}=\begin{cases}
E,&q_{kc}<0.2,\\
G,&0.2\le q_{kc}<0.4,\\
A,&q_{kc}\ge 0.4.
\end{cases}
\]

Pool B concatenates the four training snapshots and has $N$ rows.  At budget
$\rho$, a selector returns source-row indices $S$, non-negative weights $w_i$,
and a realised size $m\le\lfloor\rho N\rfloor$.  Exact-budget methods require
equality; native methods with smaller output report their realised size and are
compared at that size.  Effective learner weights are normalised to mean one.

UL is macro-F1 over enrolment labels, accompanied by MCC, balanced accuracy,
per-class precision/recall/F1, PR-AUC, log loss, and calibration.  Course
quality is estimated by probabilistic class counts (PCC):
\[
\widehat Q_c=\sum_{j\in\{E,G,A\}}\widehat\pi_{cj}\,\mu_j,
\]
where $\mu_j$ is computed from training enrolments only.  EC is the mean
absolute difference between $\widehat Q_c$ and the observed course mean among
courses satisfying the approved minimum-enrolment threshold.  Class-count
aggregation is retained as a sensitivity analysis.

\section{Data and Protocol}
\subsection{CQ V3 and Phase Snapshots}
CQ V3 is derived from a fixed MOOCCubeX project snapshot.  It removes 129,400
labelled learner--course pairs whose only observed chapter is the unresolved
bucket, retaining 178,598 labelled pairs, 168,410 learners, and 1,808 courses.
The primary target uses $\alpha=1$ in the registered quality formula.  P1--P4
use cumulative windows ending at 25\%, 50\%, 75\%, and 90\% of the observed
learner--course interval.  Labels remain fixed while features are recomputed at
each cutoff.  The target is formula-derived and has not been externally
validated as expert-rated course quality.

The M0 bundle is private and physically excludes test Parquet files.  It also
excludes full-cohort feature shards, sequences, intermediate stages, and the
overlapping default split.  A manifest records a SHA-256 digest for every
included file.
\subsection{Splits and Leakage Controls}
The selected source split assigns 1,265/181/362 courses and
128,407/16,301/33,890 learner--course pairs to train/dev/test.  No course ID is
shared across splits.  This is a random course split, not a temporal or SGKF
split; its use beyond conditional M0 evidence requires supervisor approval.
Whether \texttt{course\_id} is identical to the intended
\texttt{course\_root\_id} also remains an explicit approval item.  Test is
sealed until an approved M1 freeze.

The learner view contains 57 numeric and five categorical variables.  The
selector view contains 56 numeric behavioural variables: it removes
\texttt{phase\_fraction}, all identifier codes, target columns, and
target-derived quantities.  Median imputation and standardisation are fitted
once on the concatenated P1--P4 training pool and then reused for every dev
phase.  Missingness masks are diagnostic only and never enter either production
view.
\subsection{Endpoints and Protocol Constants}
Class-conditional means $\mu_E,\mu_G,\mu_A$ are estimated from training data and
must satisfy both their registered class intervals and
$\mu_E<\mu_G<\mu_A$.  The UL non-inferiority margin is fixed at 0.005.  The
minimum course size and EC margin remain provisional until the V3 data card and
oracle floor are approved; we do not tune them on test.

\section{Method}
\todoitem{Describe only the final approved candidate. Keep unselected candidates as ablations.}
\subsection{Selection and Allocation}
\todoitem{Mathematical objective, allocation constraints, exact/realized budget, and weight normalization.}
\subsection{Complexity and Resource Policy}
\todoitem{Measured 50k/200k/800k calibration, log--log extrapolation, and three-hour selector cap.}

\section{Experimental Design}
M0 compares FullTrain-B under natural and balanced class weighting using
logistic regression and XGBoost.  The headroom control is stratified random
selection at 0.5\%, 1\%, 2\%, and 5\%, with selector seeds 11, 29, and 47.
The M1 panel includes uniform and phase--class random controls, Gonzalez
k-center, adapted CoreTab-XGB OVA, BDIS at realised size, Moderate-DS, CCS,
Data-OOB, and classwise Kernel Thinning.  Blocked, failed, timeout, and OOM runs
remain rows in the result table and are never replaced by proxies.

The paired development unit is dataset--phase--selector seed--learner; phases
are repeated evaluations of one fitted model, not independent replicates.
Confirmatory analysis will report paired differences, clustered bootstrap
confidence intervals by course, Friedman tests across the frozen panel, and
pre-registered Wilcoxon contrasts with Holm correction.  These analyses remain
disabled until M1 freeze.

\section{Results}
\subsection{Development Results}
\paragraph{Migration to CQ V3 (21 October 2026).}
The current M0 pipeline uses a 104,580,461-byte development bundle with SHA-256
\nolinkurl{9F60946C140B7E16EE1F9BB3D94FC21BBEC898F476BE669416EBD6E4C55995A6}.
Sixty-one local protocol tests passed.  The course-ID-only P4 probe matched the
majority macro-F1 (0.269207).  A balanced multinomial missing-mask probe, used
for diagnosis only, obtained mean macro one-vs-rest ROC-AUC 0.610703 across the
four dev phases.  A private cutoff checklist sampled five cases per phase;
all 20 had maximum feature-event time no later than cutoff and zero recorded
out-of-window events.  Shuffled-label XGBoost, V3 FullTrain-B, and the V3
headroom decision are executed by the pinned Kaggle T4 cell and are not filled
from the V2 numbers below.

\paragraph{M0 headroom evidence (21 October 2026).}
The registered development-only headroom run compared FullTrain-B with C02 at
four training budgets.  XGBoost with natural class weighting was reused across
the four temporal phases; selector seeds were 11, 29, and 47.  All 52 expected
evaluation records completed, all 54 protocol tests passed, and no test metric
was opened.  Table~\ref{tab:m0-headroom} reports the phase-aggregated
development endpoint used by the registered decision rule.

\begin{table}[t]
\centering
\caption{M0 development-only headroom evidence.  The reference FullTrain-B UL
is 0.654338.  The gap is FullTrain-B UL minus mean C02 UL; none of the four
budgets reaches the registered 0.01 threshold.}
\label{tab:m0-headroom}
\begin{tabular}{rrrr}
\toprule
Budget $\rho$ & Mean C02 UL & Gap & Eligible \\
\midrule
0.5\% & 0.650153 &  0.004185 & no \\
1\%   & 0.655179 & -0.000842 & no \\
2\%   & 0.647283 &  0.007055 & no \\
5\%   & 0.652001 &  0.002337 & no \\
\bottomrule
\end{tabular}
\end{table}

Because no candidate budget met the threshold, the preregistered fallback rule
selected the smallest budget, $\rho^\star=0.005$.  This is a protocol decision,
not evidence that 0.5\% dominates every larger budget: the three selector seeds
show material variability.  The earlier 2\% screening remains diagnostic only
and cannot be used as the QD16 resource decision at $\rho^\star$.

\paragraph{Evidence integrity and current status.}
The immutable archive is
\nolinkurl{skq_cq_v1_m0_headroom_2026_10_21.tar.gz}, with SHA-256
\nolinkurl{2839B7D1AF1284999DD1AC1ED406058A9D875E7048A5FB77E2C82AD3EE4EE4A3},
produced at repository commit
\nolinkurl{9b07f6f769d8504a2555bf6608411c8f5131f8be}.  M0 remains
\textbf{PARTIAL}: QD12 headroom is complete, while QD16 still requires the ten
fixed selectors and three candidates at $\rho^\star=0.005$; FullTrain-A, LR
tuning, group-unit confirmation, and the remaining raw-data checks are also
open.  Therefore these data do not authorize M1 freezing or test opening.

\paragraph{Partial QD16 checkpoint.}
At $\rho^\star=0.005$, nine of the thirteen registered selectors completed all
three calibration sizes and three replicates.  All 17 QD16 tests passed.  The
largest completed upper runtime predictions were 5216.48 seconds for adapted
CoreTab-XGB OVA, 2128.58 seconds for GCoreSet, and 593.04 seconds for Data-OOB;
Data-OOB had an upper predicted peak RSS of 20.82 GiB.  Kernel Thinning and the
three SKQ candidates remain unresolved, so no complete development/test
resource schedule or QD16 PASS is reported.
\subsection{Confirmatory Results}
\todoitem{Populate only after M1 freeze and one-time test opening.}
\subsection{Ablations and Failure Analysis}
\todoitem{Rare-class recall, class-mass error, QP versus uniform, and runtime/quality Pareto.}

\section{Discussion}
The protocol separates evidence quality from model quality.  A successful
notebook does not imply M0 GO if the split, group unit, or resource panel is
unresolved.  The present target is an operational proxy, the split is random by
course, users may recur across splits, and course end is defined retrospectively
from observed activity.  Missingness is predictive, the rare class is below one
percent in the full V3 cohort, and course-level EC depends strongly on the
approved minimum course size.  These limitations preclude claims of causal,
cross-platform, or expert-rated course quality.

\section{Conclusion}
M0 establishes an auditable development protocol and identifies the decisions
required before confirmatory evaluation.  The next stage may proceed on
train/dev only after conditional approval; test remains sealed until the group
unit, split deviation, FullTrain-A treatment, resource decision, and M1 freeze
are resolved.

\bibliographystyle{splncs04}
\bibliography{references}
\end{document}
